\usepackage{xcolor}
\usepackage{booktabs}
\usepackage{multirow}
\usepackage[inline]{enumitem}
\usepackage{color}
\usepackage{adjustbox}
\usepackage{array}
\usepackage{tabularx} 
\usepackage{xspace}
\usepackage{etoolbox} % for acronym first-use toggles
% \usepackage{amssymb}
\usepackage{amsmath}
\usepackage{booktabs}
\usepackage{float}
\usepackage{hyperref}
\usepackage{tcolorbox}
\tcbuselibrary{skins,breakable}
\hypersetup{
    colorlinks=true,
    linkcolor=blue,
    urlcolor=cyan,
    citecolor=blue,
    anchorcolor=black,
    pdfborder={0 0 0},
}

\usepackage{graphicx}%
\usepackage{wrapfig} % enable wrapfigure/wraptable for half-width floats with text
\usepackage{subcaption}
\usepackage[normalem]{ulem}
% \usepackage{tabularray} % Not available in TeX Live 2019
% Handle newtxmath and amssymb conflict
\let\Bbbk\relax % Clear any existing definition
\usepackage{amssymb} % for checkmark and xmark
\usepackage{pifont} % for ding symbols
% \newcommand{\checkmark}{\ding{51}} % Already defined
\newcommand{\xmark}{\ding{55}}
\usepackage{booktabs}
\usepackage{subcaption}
\usepackage{listings} 
\usepackage{lstautogobble}
% 首次：使用 finalizecache 生成缓存（需 -shell-escape）；随后可切换为 frozencache
% 生成缓存后，改为 frozencache：无需 shell-escape 亦可稳定复编
\usepackage[frozencache,cachedir=_minted-main]{minted}
\usemintedstyle{vs}  % Visual Studio style for code highlighting
\setminted{
    frame=single,
    framesep=4pt,
    framerule=0.4pt,
    xrightmargin=10pt,
    breaklines=false,
    fontsize=\small,
    baselinestretch=1.0,
    escapeinside=||,
    numbers=none,
    autogobble,
    fontfamily=tt,
    samepage=true
}
\usepackage{tikz}
\usepackage{caption}
\usepackage{pict2e}
% Compact global spacing for figures/tables and captions
% Keep conservative values to avoid collisions with acmart
\captionsetup{skip=2pt,belowskip=0pt}
\captionsetup[sub]{skip=2pt}
\setlength{\abovecaptionskip}{2pt plus 1pt minus 1pt}
\setlength{\belowcaptionskip}{2pt plus 1pt minus 1pt}
\setlength{\textfloatsep}{8pt plus 2pt minus 2pt}
\setlength{\intextsep}{6pt plus 2pt minus 2pt}
\setlength{\floatsep}{6pt plus 2pt minus 2pt}
\setlength{\dbltextfloatsep}{8pt plus 2pt minus 2pt}
\setlength{\dblfloatsep}{6pt plus 2pt minus 2pt}
%\usepackage[linesnumbered,ruled,vlined]{algorithm2e}
\usepackage{algorithm}
\usepackage{multirow}
\usepackage{mathtools}	
\usepackage{mdframed}
\usepackage{syntax}
% \usepackage[skins]{tcolorbox}
\usepackage{pgfplots}
\pgfplotsset{
	compat=1.4,
	small,
	legend style={
		at={(0.99,0.99)},
		anchor=north east,
		font=\bfseries,
	},
	label style={font=\sffamily\small\bfseries}
}
\usepackage[normalem]{ulem}
\let\UnderScore_

\newcommand\redout{\bgroup\markoverwith
	{\textcolor{red}{\rule[.45ex]{1.5pt}{1.pt}}}\ULon}

\global\mdfdefinestyle{default}{%
	usetwoside=false,%
	linecolor=black,linewidth=1pt,%
	innerleftmargin=5pt,innerrightmargin=5,%
%	leftmargin=-120pt,rightmargin=-120pt,%
	everyline=true
}

\newcommand\Myperm[2][^n]{\prescript{#1\mkern-2.5mu}{}P_{#2}}
\newcommand\Mycomb[2][^n]{\prescript{#1\mkern-0.5mu}{}C_{#2}}


\newcommand{\tabincell}[2]{\begin{tabular}{@{}#1@{}}#2\end{tabular}}

\theoremstyle{definition}
\newtheorem{definition}{Definition}[section]
\newtheorem{problem}{Problem}[section]
\newtheorem{theorem}{Theorem}[section]
\newtheorem{corollary}{Corollary}[section]

\newcommand\mycommfont[1]{\small\ttfamily\textcolor{blue}{#1}}
%\SetCommentSty{mycommfont}

\makeatletter
\newenvironment{btHighlight}[1][]
{\begingroup\tikzset{bt@Highlight@par/.style={#1}}\begin{lrbox}{\@tempboxa}}
	{\end{lrbox}\bt@HL@box[bt@Highlight@par]{\@tempboxa}\endgroup}

\newcommand\btHL[1][]{%
	\begin{btHighlight}[#1]\bgroup\aftergroup\bt@HL@endenv%
	}
	\def\bt@HL@endenv{%
	\end{btHighlight}%   
	\egroup
}
\newcommand{\bt@HL@box}[2][]{%
	\tikz[#1]{%
		\pgfpathrectangle{\pgfpoint{1pt}{0pt}}{\pgfpoint{\wd #2}{\ht #2}}%
		\pgfusepath{use as bounding box}%
		\node[anchor=base west, fill=orange!25,outer sep=.5pt,inner xsep=0.5pt, inner ysep=0.15pt, rounded corners=1pt, minimum height=\ht\strutbox-.1pt,#1]{\raisebox{.01pt}{\strut}\strut\usebox{#2}};
	}%
}
\makeatother

\definecolor{add-line}{rgb}{0.84, 0.89, 0.97}
\definecolor{add-word}{rgb}{0.5, 0.7, 0.98}
\definecolor{ForestGreen}{RGB}{63,147,88}
\definecolor{remove-line}{rgb}{1.0, 0.93, 0.73}
\definecolor{remove-word}{rgb}{1.0, 0.8, 0.2}
\definecolor{Gray}{gray}{0.9}
%\definecolor{git-remove}{rgb}{red!20}
%\definecolor{git-add}{rgb}{green!20}}
\definecolor{orange}{rgb}{0.84, 0.89, 0.97}

% inline author comment macro (set to empty to hide in camera-ready)
\newcommand{\lz}[1]{\textcolor{ForestGreen}{#1}}

\newcommand{\lstbg}[3][0pt]{{\fboxsep#1\colorbox{#2}{\strut #3}}}

\lstdefinestyle{mystyle}{  
	% backgroundcolor=\color{backcolour},   
	frame=single, 
    % basicstyle=\tiny\ttfamily\bfseries,
	framexleftmargin=0pt,
	commentstyle=\color{ForestGreen},
	keywordstyle=\color{blue}\bfseries,
	numberstyle=\normalsize\color{gray},
	stringstyle=\color{purple},
	basicstyle=\normalsize\ttfamily\bfseries,
	breakatwhitespace=false,         
	breaklines=false,                 
	captionpos=b,                    
	keepspaces=true,     
	numbers=none,                    
	numbersep=4pt,               
	showspaces=false,                
	showstringspaces=false,
	showtabs=false,                  
	tabsize=2,  
	language=C++,  
	escapechar=\@,  
	%moredelim=**[is][{\btHL[fill=light-gray]}]]{#}{#},
	%moredelim=**[is][{\btHL[fill=chromeyellow]}]]{'}{'},
	%morecomment=[f][\lstbg{light-yellow}]-,
    %morecomment=[f][\lstbg{light-gray}]+,
	moredelim=**[is][{\btHL[fill=red!40]}]{@}{@},%,draw=red,dashed,thick	
	moredelim=**[is][{\btHL[fill=remove-line]}]{|*}{*|},
	moredelim=**[is][{\btHL[fill=remove-word]}]{|^}{^|},
	moredelim=**[is][{\btHL[fill=add-line]}]{||^}{^||},	
	moredelim=**[is][{\btHL[fill=add-word]}]{||*}{*||}
    % moredelim=**[is][{\btHL[fill=orange]}]{#}{#}
} 
\let\origthelstnumber\thelstnumber
\makeatletter
\newcommand*\Suppressnumber{%
  \lst@AddToHook{OnNewLine}{%
    \let\thelstnumber\relax%
     \advance\c@lstnumber-\@ne\relax%
    }%
}

\lstdefinestyle{inlinestyle}{
	language=C++,
    basicstyle=\small\ttfamily\bfseries,
}

\newcommand*\Reactivatenumber[1]{%
  \setcounter{lstnumber}{\numexpr#1-1\relax}
  \lst@AddToHook{OnNewLine}{%
   \let\thelstnumber\origthelstnumber%
   \refstepcounter{lstnumber}%
  }%
}
\makeatother
\DeclareMathOperator*{\argmax}{argmax}
\DeclareMathOperator*{\argmin}{argmin}

%\usepackage{algorithmic}
\usepackage{algpseudocode}
\renewcommand*\Call[2]{\textproc{#1}(#2)}
%\newcommand{\LineComment}[1]{\Statex \hfill /* #1 */}
\algnewcommand{\LineComment}[1]{\Statex #1}
%\algnewcommand{\LineComment}[1]{\Statex \hskip\ALG@defaultindent \(\triangleright\) #1}
\algrenewcommand\algorithmicindent{1.0em}
\algdef{SE}[DOWHILE]{Do}{doWhile}{\algorithmicdo}[1]{\algorithmicwhile\ #1}%


\newcommand*\circled[1]{\tikz[baseline=(char.base)]{
            \node[shape=circle,draw,inner sep=0.1pt] (char) {#1};}}

\usetikzlibrary{shapes,positioning,arrows.meta, calc, decorations.pathreplacing, matrix}
\usepackage{colortbl}
\definecolor{Gray}{gray}{0.85}
\tikzset{draw half paths/.style 2 args={%
  decoration={show path construction,
    lineto code={
      \draw [#1] (\tikzinputsegmentfirst) -- 
         ($(\tikzinputsegmentfirst)!0.5!(\tikzinputsegmentlast)$);
      \draw [#2] ($(\tikzinputsegmentfirst)!0.5!(\tikzinputsegmentlast)$)
        -- (\tikzinputsegmentlast);
    }
  }, decorate
}}
\usepackage{tikz-qtree}


%======font
%\usepackage{inconsolata}
%\usepackage{fontspec}
%\renewcommand*{\ttdefault}{Inconsolata}
%\let\OldTexttt\texttt
%\renewcommand{\texttt}[1]{{{\bfseries\itshape #1}}}

%\DeclareTextFontCommand{\textmtt}{\inconsolata}
\newcommand{\textmtta}[1]{{\fontfamily{zi4}\selectfont #1}}
\newcommand{\textmttb}[1]{{\fontfamily{ul9}\selectfont #1}}
\newcommand{\textmttc}[1]{{\fontfamily{jkptt}\selectfont #1}}
\newcommand{\textmttd}[1]{{\fontfamily{fvs}\selectfont #1}} %old fvm
%\newcommand{\textmtte}[1]{{\small\fontfamily{txtt}\selectfont #1}}
\newcommand{\textmtte}[1]{{\fontsize{9}{11}\fontfamily{txtt}\selectfont #1}}
\newcommand{\textmttf}[1]{{\fontfamily{ztm}\selectfont #1}}
\newcommand{\textmttg}[1]{{\fontfamily{cmvtt}\selectfont #1}}

%\renewcommand*\ttdefault{txtt}
%\renewcommand*\familydefault{\ttdefault} %% Only if the base font of the document is to be typewriter style
%\usepackage[T1]{fontenc}


\renewcommand{\texttt}[1]{\textmtte{#1}}
% Note: use \myinline|foo{}| or \myinline{a || b}.
\newcommand{\myinline}{\lstinline[style=inlineStyle]}
% \def\myinline{\lstinline[style=inlineStyle]}
\DeclareTextFontCommand{\mytexttt}{\small\ttfamily\bfseries}


%\newcommand*{\textmttaa}{\fontfamily{zi4}\selectfont}
%
%\newcommand*{\textmttbb}{\fontfamily{ul9}\selectfont}
%
%\newcommand*{\textmttcc}{\fontfamily{jkptt}\selectfont}
%
%\DeclareTextFontCommand{\textmtta}{\textmttaa}
%\DeclareTextFontCommand{\textmttb}{\textmttbb} 
%\DeclareTextFontCommand{\textmttc}{\textmttcc}

\newcommand{\myparagraph}[1]{\vspace{3pt}
\noindent
\textbf{\textit{#1}}\, 
}